\chapter{Entorno y reproducibilidad}
\label{anx:reproducibilidad}

\section{Servicios}

Todos los servicios se despliegan con Docker Compose sobre un único anfitrión y se agrupan en
perfiles, de modo que solo se levanta lo necesario para cada tarea.

\begin{table}[H]
\centering
\small
\begin{tabular}{|l|l|l|l|}
\hline
\textbf{Servicio} & \textbf{Versión} & \textbf{Perfil} & \textbf{Estado} \\
\hline
Redpanda y su consola & 25.2.1 / 3.1.0 & \texttt{core} & En uso \\
Apache Flink (gestor y trabajador) & 2.2.1 & \texttt{stream} & En uso \\
ClickHouse & 25.3 & \texttt{store} & En uso \\
Prometheus y Grafana & 3.1.0 / 11.5.1 & \texttt{observability} & En uso \\
MLflow & 2.19.0 & \texttt{mlops} & En uso \\
Portal web & --- & \texttt{portal} & En uso (demostración) \\
Neo4j con GDS & --- & \texttt{graph} & Fase 5 \\
OpenCTI & --- & \texttt{cti} & Fase 4 \\
Dagster & --- & \texttt{orchestration} & Fase 4 \\
FastAPI & --- & \texttt{api} & Fase 6 \\
\hline
\end{tabular}
\caption{Servicios del sistema y perfiles de despliegue.}
\label{tab:servicios}
\end{table}

\section{Topics del bus de eventos}

\begin{table}[H]
\centering
\small
\begin{tabular}{|l|c|l|}
\hline
\textbf{Topic} & \textbf{Particiones} & \textbf{Contenido} \\
\hline
\texttt{raw.host.windows} & 3 & Telemetría cruda de los equipos \\
\texttt{norm.events} & 6 & Eventos normalizados a OCSF \\
\texttt{det.alerts} & 3 & Alertas de las reglas Sigma \\
\texttt{det.p1\_flujo}, \texttt{det.p2\_flujo}, \texttt{det.p3\_flujo} & 3 & Evidencia de cada patrón, en orden de tiempo \\
\texttt{det.campaigns} & 3 & Campañas en el esquema común \\
\texttt{dlq.*} & 1 & Rechazos de cada etapa, con su motivo \\
\hline
\end{tabular}
\caption{Topics empleados en el Pilar A.}
\end{table}

La partición de \texttt{norm.events} se decide por la identidad del evento (conjunto de
datos, equipo, canal y número de registro), que reparte la carga de forma uniforme entre los
detectores en paralelo. Los \textit{topics} del corpus se crean sin caducidad: la retención
por defecto de siete días los vació en una ocasión sin producir ningún error.

\section{Reproducción del corpus}

El corpus de evaluación se carga con un único guion, que recrea los \textit{topics},
reproduce los nueve conjuntos de datos y los procesa con grupos de consumo nuevos:

\begin{lstlisting}[language=bash,caption={Carga del corpus del Pilar A.},label={lst:carga}]
docker compose --profile core --profile stream --profile store up -d
DETECT_INSTANCIAS=6 ./scripts/load_f3_corpus.sh
\end{lstlisting}

Las cifras esperadas tras la carga son las de la Tabla~\ref{tab:corpus-f3}. Como los
identificadores de eventos, alertas y campañas son deterministas y las tablas del almacén
deduplican por ellos, la carga puede repetirse o reanudarse sin duplicar resultados.
